\documentclass[11pt,a4paper]{article}

% --- Packages -------------------------------------------------------------
\usepackage[utf8]{inputenc}
\usepackage[T1]{fontenc}
\usepackage{lmodern}
\usepackage{amsmath,amssymb,amsthm}
\usepackage{mathtools}
\usepackage{microtype}
\usepackage[hidelinks]{hyperref}
\usepackage{url}
\usepackage{xspace}

% --- Theorem environments -------------------------------------------------
\theoremstyle{plain}
\newtheorem{theorem}{Theorem}[section]
\newtheorem{proposition}[theorem]{Proposition}
\newtheorem{lemma}[theorem]{Lemma}
\newtheorem{corollary}[theorem]{Corollary}

\theoremstyle{definition}
\newtheorem{definition}[theorem]{Definition}

\theoremstyle{remark}
\newtheorem{remark}[theorem]{Remark}

% --- Project-specific macros ----------------------------------------------
\newcommand{\Cl}{\mathrm{Cl}}
\newcommand{\Spin}{\mathrm{Spin}}
\newcommand{\Pin}{\mathrm{Pin}}
\newcommand{\SO}{\mathrm{SO}}
\newcommand{\End}{\mathrm{End}}
\newcommand{\Mat}{\mathrm{Mat}}
\newcommand{\dualProd}{\mathrm{dualProd}}
\newcommand{\wedgeW}{\bigwedge W}
\newcommand{\wedgeEven}{\bigwedge^{\mathrm{even}} W}
\newcommand{\wedgeOdd}{\bigwedge^{\mathrm{odd}} W}
\newcommand{\SpinorLean}{\textsc{SpinorLean}\xspace}
\newcommand{\Lean}{Lean~4\xspace}
\newcommand{\Mathlib}{\textsc{Mathlib}\xspace}
\newcommand{\code}[1]{\texttt{#1}}

% --- Metadata -------------------------------------------------------------
\title{Formalizing Spinor Representations from Clifford Algebras in
       \Lean\ and \Mathlib}
\author{%
  Author One\\\small Affiliation One\\\small \texttt{author.one@example.org}
  \and
  Author Two\\\small Affiliation Two\\\small \texttt{author.two@example.org}
}
\date{\today}

\begin{document}
\maketitle

% --- Abstract -------------------------------------------------------------
\begin{abstract}
We present \SpinorLean, the first formalization of spinor representations
from Clifford algebras in a proof assistant. Building on \Mathlib's existing
\code{CliffordAlgebra}, \code{ExteriorAlgebra}, and \code{spinGroup}
libraries, we construct the spinor module as the exterior algebra
$\wedgeW$ of a maximal totally isotropic subspace $W \le V$, equipped with
its natural $\Cl(V,Q)$-action via wedge and contraction. The formalization
is organized around a hyperbolic-presentation pattern: an explicit isometry
$Q \simeq \dualProd_K W$ is packaged as first-class data that transports the
chosen $\wedgeW$ model, its Clifford action, its $\Spin(Q)$-action, its
parity splitting, and its simplicity and non-factorization results from the
concrete split model to arbitrary quadratic forms that admit such a
presentation. The resulting library includes faithful Clifford actions,
half-spin (Weyl) simplicity and inequivalence, the algebra equivalences
$\Cl(Q) \simeq \End(\wedgeW)$ and $\Cl^+(Q) \simeq
\End(\wedgeEven) \times \End(\wedgeOdd)$, matrix models over $\mathbb{C}$
and split real signatures (including $\Cl(n,n) \simeq \Mat_{2^n}(\mathbb{R})$),
explicit low-dimensional identifications, and the compact identification
$\Spin(2) \simeq U(1)$. The library builds cleanly with zero
\code{sorry}/\code{admit}.
\end{abstract}

% --- 1. Introduction -----------------------------------------------------
\section{Introduction}
\label{sec:intro}

Spinors sit at the centre of twentieth-century mathematics and physics.
Algebraically, they are the ``square roots of geometry'': the irreducible
modules over a Clifford algebra $\Cl(V,Q)$ on which the spin group
$\Spin(V,Q)$ acts. Through Atiyah--Bott--Shapiro, they power real
$K$-theory and Bott periodicity; through Dirac operators, they describe
fermionic matter fields; through twistors and pure spinors, they appear
throughout modern geometry. Any foundational library of geometry in a proof
assistant eventually needs them.

Despite this, no proof assistant has previously formalized the spinor
representation in the precise sense: a $\Cl(V,Q)$-module $S$ carrying, by
restriction, a representation of $\Spin(V,Q)$, with dimension $2^{n/2}$
and the expected simplicity and chirality properties. \Mathlib~\cite{Mathlib}
already contains substantial Clifford-algebra infrastructure --
\code{CliffordAlgebra}, \code{ExteriorAlgebra}, \code{spinGroup},
\code{pinGroup}, and (via Wieser~\cite{Wieser2022Clifford}) their universal
properties and gradings -- but the Clifford \emph{module} that connects
$\Cl(V,Q)$ and $\Spin(V,Q)$ to linear algebra was missing. The earlier
\code{lean-ga} library~\cite{LeanGA} concentrated on computational geometric
algebra in \Lean~3 and did not reach the representation-theoretic
content.

\SpinorLean fills this gap. Given a quadratic space $(V,Q)$ over a field
$K$ of characteristic not two, we construct the exterior model
$\wedgeW$ attached to a maximal totally isotropic subspace $W \le V$,
equip it with a $\Cl(V,Q)$-action, restrict to a
$\Spin(V,Q)$-representation, split off the chiral (Weyl) halves, and prove
the corresponding simplicity, faithfulness, inequivalence, and
matrix-classification theorems. In finite-dimensional split rank the
representation has dimension $2^{n/2}$ and the even and odd halves each
have dimension $2^{n/2-1}$. All results are machine-checked in \Lean,
with a clean \code{lake build} and zero \code{sorry}/\code{admit}.

The main conceptual contribution is architectural rather than
mathematical: we isolate the \emph{hyperbolic presentation}
$Q \simeq \dualProd_K W$ as the primitive piece of data and thread a
chosen $\wedgeW$ model through it. This keeps every theorem concrete
enough to prove on the split model and generic enough to apply to any
quadratic form admitting such a presentation.

% --- 2. Related Work -----------------------------------------------------
\section{Related Work}
\label{sec:related}

\paragraph{Mathlib.} The Clifford-algebra infrastructure we build on is
extensive. \Mathlib provides \code{CliffordAlgebra Q} with its universal
property, the exterior-algebra equivalence
\code{CliffordAlgebra.equivExterior} in characteristic not two, the even
subalgebra and $\mathbb{Z}/2$-grading via \code{evenOdd}, and the groups
\code{spinGroup Q} and \code{pinGroup Q} together with their conjugation
actions. Wieser's work on the Clifford algebra
design~\cite{Wieser2022Clifford} underlies much of this API. What was
missing -- and what \SpinorLean adds -- is the Clifford \emph{module}
structure: a concrete $\Cl(V,Q)$-module realizing the spinor
representation, together with its simplicity, half-spin decomposition,
and matrix-algebra classifications.

\paragraph{lean-ga.} The \code{lean-ga} project~\cite{LeanGA}, developed
largely in \Lean~3, explored geometric algebra with a strong emphasis on
computational and basis-level constructions. It predates much of the
current \Mathlib Clifford stack and does not target spinor modules or
their representation theory; our work is complementary, focusing on the
algebraic representation-theoretic side in \Lean~4 over the modern
\Mathlib API.

\paragraph{Mathematics background.} Our development follows the classical
presentations of Lawson--Michelsohn~\cite{LawsonMichelsohn1989},
Chevalley~\cite{Chevalley1996}, Atiyah--Bott--Shapiro~\cite{AtiyahBottShapiro1964},
and Lounesto~\cite{Lounesto2001}, with the low-dimensional octonionic
and quaternionic incidences drawn from Baez~\cite{Baez2002}. The
hyperbolic-presentation pattern can be seen as a proof-engineering
refinement of the standard ``choose a maximal isotropic and an exterior
model'' construction in those references.

% --- 3. Formalization Architecture ---------------------------------------
\section{Formalization Architecture}
\label{sec:arch}

\paragraph{File layout.} The library is organized into a small number of
\code{Spinor/\ldots} files, each with a single responsibility: Mathlib
imports and notation; isotropic subspaces and Witt decomposition
(\code{Isotropic.lean}, \code{WittDecomp.lean}); the chosen exterior model
and its wedge/contraction operators (\code{ExteriorModel.lean}); the
concrete split Clifford and spin actions (\code{HyperbolicAction.lean}); the
first-class presentation API (\code{Presentation.lean}); the transport of
all Clifford/spin structure onto the chosen exterior model
(\code{CliffordAction.lean}, \code{SpinRep.lean}, \code{Chiral.lean}); the
orthogonal/covering-map layer (\code{OrthogonalAction.lean}); and the
doubled form (\code{ProdNeg.lean}). Classification results live in
\code{ComplexClassification.lean}, \code{RealClassification.lean}, and
\code{LowDimensional.lean}.

\paragraph{The hyperbolic-presentation pattern.} The central organizing
idea is to isolate the data
\[
  Q \simeq \dualProd_K W
  \qquad (W \text{ a maximal totally isotropic subspace})
\]
as a \Lean{} structure,
\code{HyperbolicPresentation Q}, carrying a subspace $W$ and an isometry
equivalence to the standard hyperbolic form \code{dualProd K W}. This
structure is promoted to a first-class object: it carries a chosen
$\wedgeW$ model (\code{spinorModule}), chosen even and odd halves
(\code{evenSpinorModule}, \code{oddSpinorModule}), the transported maximal
totally isotropic subspace, the Witt-index theorem
(\code{wittIndex\_eq\_finrank}), the transported
\code{Cl(Q) \ensuremath{\to} End(\ensuremath{\wedgeW})} action, and the
restricted \code{spinGroup Q} action as actual \Lean{} \code{Module} and
\code{MulAction} instances. Theorems about the ambient $(V,Q)$ are then
proved \emph{on} this presentation and reused whenever one is available.

\paragraph{Four transport layers.} Concretely, the chosen $\wedgeW$ model
is threaded through four layers that sit above each other:
\begin{enumerate}
  \item \textbf{The split model.} All hard work happens over
        $\Cl(W^* \times W, \dualProd)$ acting on $\wedgeW$ by wedge on
        the $W$ factor and interior product by the $W^*$ factor
        (\code{HyperbolicAction.lean}). The Clifford relation is verified
        directly; the resulting map is proved to be an algebra equivalence
        $\Cl(W^* \times W, \dualProd) \simeq \End(\wedgeW)$ by explicit
        basis projectors.
  \item \textbf{Explicit hyperbolic presentations
        $Q \simeq \dualProd_K W$.} Any such presentation transports the
        split action to $\Cl(V,Q) \to \End(\wedgeW)$ and to the restricted
        $\Spin(V,Q)$-representation, together with parity, Witt-index,
        simplicity, faithfulness, and non-factorization statements
        (\code{Presentation.lean}, \code{CliffordAction.lean},
        \code{SpinRep.lean}, \code{Chiral.lean}).
  \item \textbf{The Witt presentation API.} The genuine Witt decomposition
        $V \simeq W \oplus W^* \oplus V_0$, obtained from a totally
        isotropic subspace plus a chosen complement, yields both a linear
        decomposition and a quadratic splitting
        $Q \simeq \dualProd_K W \oplus Q_0$ (\code{wittIsometryEquiv}).
        Its hyperbolic factor is exposed as
        \code{HyperbolicPresentation.canonicalWittFactor} and feeds back
        into layer~2.
  \item \textbf{The split-rank canonical model.} When $V_0$ vanishes
        (split rank), the canonical Witt subspace
        \code{Q.wittSubspace} chooses its complement internally, yielding
        \code{splitWittPresentation} and \code{splitSpinorModule}, together
        with the top-level capstone theorems (dimensions, simplicity,
        faithfulness, and non-factorization) that require no explicit
        isometry argument.
\end{enumerate}
A fifth specialized layer handles the nondegenerate doubled case
$(Q \oplus -Q)$ through the canonical diagonal isotropic subspace, via
\Mathlib's \code{QuadraticForm.toDualProd}, upgraded here to an
isometry equivalence and packaged as a reusable
\code{HyperbolicPresentation}.

\paragraph{Threading the chosen model.} The effect of this architecture is
that every theorem is proved once on the concrete $\wedgeW$ model, then
lifted verbatim through the presentation APIs. For example, full-module
simplicity is proved in \code{HyperbolicAction} for the split model and
then surfaces on \code{HyperbolicPresentation}, \code{WittPresentation},
and the top-level split-rank aliases with no additional mathematical work
-- only bookkeeping via the transported isomorphisms.

% --- 4. Main Results -----------------------------------------------------
\section{Main Results}
\label{sec:results}

The following catalog summarises the capstone theorems present in the
current library.

\paragraph{Foundations.}
\begin{itemize}
  \item Totally isotropic subspaces, existence of maximal totally
        isotropic subspaces over finite-dimensional spaces over fields,
        and \code{wittIndex} as the maximal dimension of such a subspace.
  \item The general linear Witt decomposition
        $V \simeq W \oplus W^* \oplus V_0$ from a chosen complement
        (\code{wittLinearDecompositionOfIsCompl}) and the corresponding
        quadratic splitting \code{wittIsometryEquivOfIsCompl}.
\end{itemize}

\paragraph{The chosen exterior model.}
\begin{itemize}
  \item $\dim(\wedgeW) = 2^{\,\mathrm{wittIndex}\,Q}$ and, under an
        explicit hyperbolic presentation,
        $\dim(\wedgeW) = 2^{\,\dim V / 2}$.
  \item Parity splitting $\wedgeW = \wedgeEven \oplus \wedgeOdd$, with
        explicit wedge and contraction operators that flip parity; in
        positive split rank, each half has dimension $2^{\dim V/2 - 1}$.
\end{itemize}

\paragraph{Clifford and spin actions.}
\begin{itemize}
  \item Split hyperbolic action $\Cl(W^* \times W, \dualProd) \to
        \End(\wedgeW)$ satisfying the Clifford relation.
  \item Transport to $\Cl(V,Q) \to \End(\wedgeW)$ along any explicit
        hyperbolic presentation, and its restriction to
        \code{spinGroup Q} as a genuine \code{MulAction}.
  \item Faithful Clifford action on the chosen model in the explicit
        hyperbolic, Witt, and split-Witt settings, packaged as an algebra
        equivalence $\Cl(Q) \simeq \End(\wedgeW)$.
\end{itemize}

\paragraph{Chirality / half-spin.}
\begin{itemize}
  \item Chosen-model chiral decomposition with \code{spinGroup} preserving
        each half, and identification of the zero-form chiral pieces with
        the exterior even/odd summands.
  \item $\Cl^+(Q) \simeq \End(\wedgeEven) \times \End(\wedgeOdd)$ through
        the presentation APIs; in positive split rank, the concrete
        product-of-matrices form of size $2^{\dim V/2 - 1}$.
\end{itemize}

\paragraph{Simplicity, faithfulness, and covering-map statements.}
\begin{itemize}
  \item The full chosen Clifford module $\wedgeW$ is simple in the split
        model; this is transported to the hyperbolic, Witt, and
        split-Witt presentations.
  \item Each chosen half $\wedgeEven$, $\wedgeOdd$ is simple under the
        even Clifford algebra (positive split rank for the odd half),
        and the two halves are inequivalent as modules over $\Cl^+$.
  \item The ambient vector action
        \code{spinLinearRepresentation} preserves $Q$; each spin element
        is packaged as a genuine \code{Q.IsometryEquiv Q}, and the whole
        map is a homomorphism
        \code{spinIsometryRepresentation :\,spinGroup\,Q \ensuremath{\to^{*}}
        Q.IsometryEquiv\,Q}.
  \item On the split-rank chosen-model API, the ambient
        spin-to-isometry kernel is proved to be exactly
        $\{\pm 1\}$, and the spin representation is proved not to
        factor through the ambient isometry representation whenever
        $Q$ represents $-1$ and $-1 \neq 1$.
\end{itemize}

\paragraph{Classification.}
\begin{itemize}
  \item Over $\mathbb{C}$: $\Cl(2n,\mathbb{C}) \simeq
        \Mat_{2^n}(\mathbb{C})$ and
        $\Cl(2n{+}1,\mathbb{C}) \simeq
        \Mat_{2^n}(\mathbb{C}) \times \Mat_{2^n}(\mathbb{C})$ on the
        standard sum-of-squares forms.
  \item Over $\mathbb{R}$ (split signature):
        $\Cl(n,n) \simeq \Mat_{2^n}(\mathbb{R})$ and
        $\Cl^+(n,n) \simeq \Mat_{2^{n-1}}(\mathbb{R}) \times
        \Mat_{2^{n-1}}(\mathbb{R})$.
  \item Explicit low-dimensional identifications, including
        $\Cl(1,\mathbb{C})$, $\Cl(2,\mathbb{C})$, $\Cl(3,\mathbb{C})$,
        $\Cl(4,\mathbb{C})$, $\Cl(0,1) \simeq \mathbb{C}$,
        $\Cl(0,2) \simeq \mathbb{H}$, $\Cl^+(2,0) \simeq \mathbb{C}$,
        $\Cl^+(3,0) \simeq \mathbb{H}$, $\Cl(1,1) \simeq
        \Mat_2(\mathbb{R})$, and $\Cl^+(2,2) \simeq \Mat_2(\mathbb{R})
        \times \Mat_2(\mathbb{R})$.
  \item The compact group identification
        $\Spin(2) \simeq U(1)$, packaged as
        \code{realSpin02EquivUnitaryComplex}.
\end{itemize}

% --- 5. Key Proof Highlights ---------------------------------------------
\section{Key Proof Highlights}
\label{sec:highlights}

\paragraph{Simplicity of $\wedgeEven$ under $\Cl^+$.}
Half-spin simplicity is the technical heart of the development. In the
split model $\Cl(W^* \times W, \dualProd)$ acting on $\wedgeW$ by
wedge and contraction, we exhibit explicit basis projectors witnessing
that every non-zero element of $\wedgeEven$ can be mapped to any other
by an element of the even Clifford algebra. The argument proceeds via the
basis-indexed decomposition of $\wedgeW$, pairing wedge-by-$W$ with
interior-by-$W^*$ to produce matrix units acting on the even subspace. In
positive split rank the same argument yields simplicity of $\wedgeOdd$.
This split-model statement is then transported through the
\code{HyperbolicPresentation}, \code{WittPresentation}, and
\code{splitWittPresentation} layers with no additional work.

\paragraph{Equivalence $\Cl(Q) \simeq \End(\wedgeW)$.}
Surjectivity and injectivity of
\code{splitCliffordAction :
  Cl(W^* \times W, dualProd) \ensuremath{\to} End(\ensuremath{\wedgeW})}
are the counterpart of simplicity from the algebra side. Surjectivity is
proved by writing an arbitrary endomorphism as a linear combination of
matrix units built from the same wedge/contraction projectors; injectivity
follows because the split action is faithful on a spanning set. Combining
with the transport along any hyperbolic presentation yields the algebra
equivalence $\Cl(Q) \simeq \End(\wedgeW)$ and, in split rank, the
concrete $\Cl(Q) \simeq \Mat_{2^{\dim V/2}}(K)$ realization through
\code{splitWittCliffordEquivMatrix}. The same machinery, applied to
$\Cl^+$, yields
$\Cl^+(Q) \simeq \End(\wedgeEven) \times \End(\wedgeOdd)$.

\paragraph{Split-rank covering-map fragments.}
The covering map $\Spin(V,Q) \to \SO(V,Q)$ is not yet formalized in full
generality, but its key algebraic ingredients are. We show that every
element of the kernel of the spin-to-isometry representation acts
trivially on the whole Clifford algebra by conjugation, hence commutes
with every Clifford element. Over domains this forces the scalar part of
the kernel to be $\pm 1$, and in the finite-dimensional hyperbolic setting
the chosen-model algebra equivalence $\Cl(Q) \simeq \End(\wedgeW)$
forces every kernel element to be scalar. Combining the two yields, on
the canonical split-rank API, the exact kernel
$\{\pm 1\}$ via
\code{splitSpinorCoveringKernel\_eq\_one\_or\_neg\_one} and the
non-factorization result
\code{splitSpinorRepresentation\_not\_factor\_through\_isometry}
whenever $Q$ represents $-1$ and $-1 \neq 1$.

% --- 6. Lessons Learned --------------------------------------------------
\section{Lessons Learned}
\label{sec:lessons}

\paragraph{Gaps found in \Mathlib.} The inventory in
\code{Spinor/Inventory.lean} records three concrete gaps that \SpinorLean
had to bridge. First, \code{QuadraticForm.WittDecomp} was effectively
absent: the existence of maximal totally isotropic subspaces over
finite-dimensional vector spaces, the Witt index as the maximal such
dimension, and the linear/quadratic Witt decomposition $V \simeq W \oplus
W^* \oplus V_0$ had to be built from primitive pieces and are now exposed
as \code{wittLinearDecompositionOfIsCompl} and \code{wittIsometryEquivOfIsCompl}.
Second, interior products on the exterior algebra were only partially
present: \code{ExteriorAlgebra.\ensuremath{\iota}Multi} existed, but a
clean contraction operator $\iota_f : \wedgeW \to \wedgeW$
dual to wedge multiplication, with the relations used in the Clifford
action, had to be developed in \code{ExteriorModel.lean}. Third, and most
importantly, the Clifford \emph{module} structure -- the spinor
representation itself -- was entirely missing; this is the main technical
contribution of the library.

\paragraph{Design choices.} The hyperbolic-presentation pattern emerged
from two failed drafts in which the chosen model was threaded either too
concretely (forcing every consumer to carry a subspace and isometry by
hand) or too abstractly (losing access to the explicit basis projectors
needed for simplicity proofs). Packaging the isometry and its chosen
model as a single \Lean{} structure, and duplicating exactly the API
needed at each transport layer, gave us both reusable theorems and
manageable proof obligations. A recurring theme was the benefit of
proving each hard statement once on the split model and letting the
transport layer handle the rest -- the four-layer stack has paid back
that investment many times over.

% --- 7. Future Work ------------------------------------------------------
\section{Future Work}
\label{sec:future}

Within the scope of the current library, several natural extensions
remain. First, the general \emph{covering map} $\Spin(V,Q) \to \SO(V,Q)$
is not yet formalized in full generality: the kernel $\{\pm 1\}$ is
established on the canonical split-rank API, but a fully ambient
surjectivity statement remains open. Second, \emph{Bott periodicity}
for real Clifford algebras (period 8) is still only foundational: the
split signature $\Cl(n,n)$ classification is in place, but the genuinely
eight-fold periodicity phenomenon has not been formalized. Third, the
remaining \emph{low-dimensional spin identifications}
$\Spin(3)$, $\Spin(4)$, $\Spin(5)$, and $\Spin(6)$ -- connecting spinors
to quaternions, pairs of quaternions, and symplectic/unitary groups --
would round out the explicit catalog beyond the already-packaged
$\Spin(2) \simeq U(1)$.

Beyond the current algebraic scope, the \code{ROADMAP.md} non-goals
section lists \emph{spinor bundles}, \emph{Dirac operators}, and the
associated differential-geometric content as out of scope for the
present paper, awaiting a more mature \Mathlib manifold and connection
theory. Computational paths, basis-level geometric-algebra calculations,
and direct physics applications are similarly left for follow-up work.
We view the current development as the algebraic substrate that makes
those future directions approachable, not as their replacement.

% --- Bibliography --------------------------------------------------------
\bibliographystyle{alpha}
\bibliography{refs}

\end{document}
